\section{Discussion}\label{sec:discussion}

\subsection{Why persistence does not improve forecasts}

Section~\ref{sec:window} shows that the rolling persistence state is largely
a smoothed record of the largest volatility episodes of the preceding two to
four years. It rises when such an episode enters the window and falls, on a
date fixed by the window length, when the episode leaves.
Sections~\ref{sec:results} and \ref{sec:robust} show that, once the VIX and
MOVE are in the model, the state does not significantly improve forecasts in
any of the registered designs. We interpret the first result as the
explanation for the second. Implied volatility is a forward-looking price
that already reflects recent turmoil, partly through a variance risk premium
that rises after crises and decays slowly \citep{bates2000post,
bekaert2014vix, andersen2015risk}. A backward-looking indicator of when the
last crisis occurred has little to add. The Clark--West rejections show that
it adds some information in population; the Giacomini--White tests show that
this contribution is too small to survive estimation.

These results do not imply that volatility lacks long memory, or that the
persistence of volatility shocks never changes. The full-sample estimates in
Table~\ref{tab:memory} are large, and our short-memory null reproduces their
calm-period level, consistent with a long literature showing that level
shifts and long memory are difficult to separate \citep{diebold2001long,
perron2010long, varneskov2018combining}. Our claim is narrower: the time
variation in rolling estimates is not evidence of time variation in
persistence. It is the pattern that a process with constant parameters
produces when a few large episodes pass through a fixed window.

\subsection{Implications for applied work}

Three practices follow for research that uses rolling memory estimates,
Hurst exponents or similar windowed statistics. First, results should be
reported for several window lengths, with a check of whether the turning
points move with the window; a structural change should not. Second, the
windowed statistic can be regressed on the largest value of an observable
stress measure inside the same window; a high $R^2$ indicates that the
statistic reflects window composition rather than dynamics. Third, the
series can be compared with a short-memory null into which the observed
episodes are injected, and with estimators and tests designed to separate
long memory from level shifts \citep{qu2011test, hou2014modified}. These are
standard tools in fields that encountered the same problem earlier, from
hydrology \citep{klemes1974hurst} to early-warning indicators in ecology and
climate science \citep{dakos2008slowing, ditlevsen2010tipping,
boettiger2012quantifying} and sliding-window connectivity in neuroscience
\citep{leonardi2015spurious}.

Three lessons for forecasting studies extend beyond this application. The
benchmark and the challenger must use the same information set, down to the
day: a one-day lag in the inputs of the benchmark inflated the value of
implied volatility by about a third. When a nested comparison yields a
Clark--West rejection without a lower loss, a Giacomini--White test indicates
whether the estimated model is useful. Finally, machine-learning comparisons
are informative only if every learner is tuned under the same protocol: with
default settings gradient boosting lost 12--21\% relative to HAR, against
0.2--8\% once tuned.

\subsection{Limitations}

The panel consists of current S\&P~500 members with long histories. Firms
that left the index are missing, which removes some of the most volatile
histories; this omission affects the level of volatility and of the memory
estimates more than the window mechanism, which is a property of the
estimator. The variance proxy is the daily range rather than intraday
realized variance. Its sampling noise lowers the level of $\hat d$ and of the
Hurst exponent but does not change the dates at which episodes enter and
leave the window. Returns are price returns, which matters only for the
portfolio exercise. The forecast evaluation period, 2013 to 2026, contains
one extreme episode, in 2020, and several moderate ones; the window analysis
also uses 2008. Markets without a liquid implied-volatility index are where
persistence could add the most, but in such markets the largest recent
realized variance is directly available and would be the natural benchmark.
Finally, the evidence on the window mechanism concerns memory in volatility;
a general analysis across estimators, asset classes and markets is left for
future work.

\section{Conclusions}\label{sec:conclusions}

Rolling long-memory estimates of stock volatility rise sharply in crises, and
it is tempting to read them as a measure of how persistent the market has
become. For 115 S\&P~500 stocks over 2001--2026 they do not improve
volatility forecasts beyond implied volatility under any design we
registered: per-stock and pooled regressions, the level and the term
structure of future variance, a market-level series, linear models and tuned
machine-learning estimators, calm and stressed periods, and
volatility-managed portfolios. The evidence points to a simple reason: the
rolling persistence state records which extreme episodes lie inside its
window. It jumps when an episode enters and falls when the episode leaves, at
dates that move one-for-one with the window length; a short-memory process
reproduces this pattern; and 83\% of its variation is explained by the
largest VIX inside the window. Researchers who use rolling memory estimates
as a state variable should check for these signatures before interpreting
them. Forecasting studies that compare HAR with implied volatility should
align the information sets of the two models, which here removes a third of
the apparent value of the VIX.
